\documentclass[10pt,a4paper]{amsart}
\usepackage{amsmath,amssymb,amsthm,graphicx}
% For commutative diagrams we use Xy-pic
\usepackage[all]{xy}

% We use 2cell for 2-commutative diagrams.
\xyoption{2cell}
% Theorem environments.
%
\theoremstyle{plain}
\newtheorem{theorem}[subsection]{Theorem}
\newtheorem{proposition}[subsection]{Proposition}
\newtheorem{lemma}[subsection]{Lemma}

\theoremstyle{definition}
\newtheorem{definition}[subsection]{Definition}
\newtheorem{example}[subsection]{Example}
\newtheorem{exercise}[subsection]{Exercise}
\newtheorem{situation}[subsection]{Situation}

\theoremstyle{remark}
\newtheorem{remark}[subsection]{Remark}
\newtheorem{remarks}[subsection]{Remarks}

\numberwithin{equation}{subsection}

% Macros
%
\def\lim{\mathop{\mathrm{lim}}\nolimits}
\def\colim{\mathop{\mathrm{colim}}\nolimits}
\def\Spec{\mathop{\mathrm{Spec}}}
\def\Hom{\mathop{\mathrm{Hom}}\nolimits}
\def\Ext{\mathop{\mathrm{Ext}}\nolimits}
\def\SheafHom{\mathop{\mathcal{H}\!\mathit{om}}\nolimits}
\def\SheafExt{\mathop{\mathcal{E}\!\mathit{xt}}\nolimits}
\def\Sch{\mathit{Sch}}
\def\Mor{\mathop{\mathrm{Mor}}\nolimits}
\def\Ob{\mathop{\mathrm{Ob}}\nolimits}
\def\Sh{\mathop{\mathit{Sh}}\nolimits}
\def\NL{\mathop{N\!L}\nolimits}
\def\CH{\mathop{\mathrm{CH}}\nolimits}
\def\proetale{{pro\text{-}\acute{e}tale}}
\def\etale{{\acute{e}tale}}
\def\QCoh{\mathit{QCoh}}
\def\Ker{\mathop{\mathrm{Ker}}}
\def\Im{\mathop{\mathrm{Im}}}
\def\Coker{\mathop{\mathrm{Coker}}}
\def\Coim{\mathop{\mathrm{Coim}}}

% Boxtimes
%
\DeclareMathSymbol{\boxtimes}{\mathbin}{AMSa}{"02}

%
% Macros for moduli stacks/spaces
%
\def\QCohstack{\mathcal{QC}\!\mathit{oh}}
\def\Cohstack{\mathcal{C}\!\mathit{oh}}
\def\Spacesstack{\mathcal{S}\!\mathit{paces}}
\def\Quotfunctor{\mathrm{Quot}}
\def\Hilbfunctor{\mathrm{Hilb}}
\def\Curvesstack{\mathcal{C}\!\mathit{urves}}
\def\Polarizedstack{\mathcal{P}\!\mathit{olarized}}
\def\Complexesstack{\mathcal{C}\!\mathit{omplexes}}
% \Pic is the operator that assigns to X its picard group, usage \Pic(X)
% \Picardstack_{X/B} denotes the Picard stack of X over B
% \Picardfunctor_{X/B} denotes the Picard functor of X over B
\def\Pic{\mathop{\mathrm{Pic}}\nolimits}
\def\Picardstack{\mathcal{P}\!\mathit{ic}}
\def\Picardfunctor{\mathrm{Pic}}
\def\Deformationcategory{\mathcal{D}\!\mathit{ef}}
\title{Gluing finitely many finite local extensions along the height one primes of a two-dimensional local ring}
\author{The Stacks Project Authors}
\date{stacks-project:chow.tex:local-snapshot-20261011; GNU FDL 1.2 or later}
\begin{document}
\maketitle
\noindent\textit{Editorial note.} All delivered mathematical text is the author's own, copied from the pinned Stacks Project chapter file chow.tex; no statement or proof step was rewritten, reordered or removed. The author's label lemma-perpare-key contains a misspelling; it is reproduced verbatim in the claim key and in the source anchors and is not corrected.\par\medskip
\noindent\textit{Editorial note.} The preamble is the chapter's own native preamble excerpt used for this source: the Xy-pic setup at lines 14-18 and the theorem declarations with the mathematical macros at lines 154-217 of sources/stacks-chow/preamble.tex. The quoted class line and the xr-hyper externaldocument block are not included because the assembler supplies a plain amsart wrapper; no macro was invented or redefined.\par\medskip
\noindent\textit{Editorial note.} The counted statement and proof contain no cross reference at all, so the assembled document defines every label it uses and no external reference stub is needed.\par\medskip
\noindent\textit{Allowed standard prerequisites (editorial).}\begin{itemize}
\item Noetherian local ring of dimension two with a nonzerodivisor t, its height one primes above t. Standard commutative algebra; the statement writes V(t) = \{m, q\_1, ..., q\_r\} for the primes containing t, so the q\_i are the height one primes of the two-dimensional local ring. Supplies the geometry of the situation, in particular that only finitely many primes lie above t and that the localizations A\_\{q\_i\} are one-dimensional; no part of the gluing construction is taken from it.
\item Localization is exact and A is Noetherian. Standard commutative algebra; used at proof lines 1122-1124 to conclude that M and M' are finite A-modules with M\_\{q\_i\} = M'\_\{q\_i\} = B\_i. Supplies exactness of localization and finite generation; the construction of M, M' and B stays inside the counted proof.
\end{itemize}\medskip
\begin{lemma}
\label{lemma-perpare-key}
Let $(A, \mathfrak m)$ be a $2$-dimensional Noetherian local ring.
Let $t \in \mathfrak m$ be a nonzerodivisor. Say
$V(t) = \{\mathfrak m, \mathfrak q_1, \ldots, \mathfrak q_r\}$.
Let $A_{\mathfrak q_i} \subset B_i$ be a finite ring
extension with $B_i/A_{\mathfrak q_i}$ annihilated by a power of
$t$. Then there exists a finite extension $A \subset B$ of
local rings identifying residue fields
with $B_i \cong B_{\mathfrak q_i}$ and $B/A$ annihilated
by a power of $t$.
\end{lemma}

\begin{proof}
Choose $n > 0$ such that $B_i \subset t^{-n}A_{\mathfrak q_i}$.
Let $M \subset t^{-n}A$, resp.\ $M' \subset t^{-2n}A$ be the
$A$-submodule consisting of elements mapping to $B_i$ in
$t^{-n}A_{\mathfrak q_i}$, resp.\ $t^{-2n}A_{\mathfrak q_i}$.
Then $M \subset M'$ are finite $A$-modules as $A$ is Noetherian
and $M_{\mathfrak q_i} = M'_{\mathfrak q_i} = B_i$ as localization
is exact. Thus $M'/M$ is annihilated by $\mathfrak m^c$ for some
$c > 0$. Observe that $M \cdot M \subset M'$ under the multiplication
$t^{-n}A \times t^{-n}A \to t^{-2n}A$. Hence
$B = A + \mathfrak m^{c + 1}M$ is a finite $A$-algebra with the correct
localizations. We omit the verification that $B$ is local with
maximal ideal $\mathfrak m + \mathfrak m^{c + 1}M$.
\end{proof}

\end{document}
